\documentclass[UTF8,a4paper,12pt]{ctexart}

% 宏包
\usepackage[demo]{graphicx}  % demo选项使图片显示为占位框，无需实际图片文件
\usepackage{amsmath,amssymb,amsthm}
\usepackage{booktabs}
\usepackage{tabularx}
\usepackage{array}
\usepackage{geometry}
\usepackage{hyperref}
\usepackage{caption}
\usepackage{float}
\usepackage{enumitem}

\geometry{left=2.5cm,right=2.5cm,top=2.5cm,bottom=2.5cm}

\title{问题三：基于问题二前向机制的视觉认知动态宏观模型}
\author{参赛队}
\date{2026年9月26日}

\begin{document}

\maketitle

\begin{abstract}
题目要求在问题二脑电形成机制基础上构造认知宏观模型，并用实测脑电验证，同时指出应答可能错误或未及时。本文以四份256Hz记录中的F3、Fz、F4和VisCue为观测对象，沿“前向机制衔接、动态状态构造、整份记录留出检验”建立问题三模型。

\textbf{事件审计方面}，400个提示事件各对应一个第9通道起始沿，其相对提示时间中位数为2.2148s，接近候选目标时刻。经核实，第9通道编码为应答时刻（-2=点击左，+2=点击右），但无独立目标起始标记。因此，本文将第9通道用于确定应答时刻 \(t_{act}\) 和截取应答前100ms窗口，同时生成正确/错误/未及时标签用于扩展验证，不参与EEG特征计算和认知状态模型输入。按边界、非有限值、原始幅度和通道平直规则筛选后保留355个完整试次。

\textbf{模型方面}，本文以问题二前向链在候选提示与目标场景下生成五个源代理，经 \(3\times 5\) 导联矩阵正向投影到三电极，建立视觉驱动基线。在此基础上，引入一个具有保持、衰减和交互作用的记忆相关动态状态 \(H_k\)，通过明确的观测方程生成三通道EEG。该状态仅解释为候选认知过程，不直接等同海马活动；三个观测通道不意味着必须设置三个认知状态。

\textbf{验证方面}，每折完整留出一份记录，用其余三份记录拟合，在固定提示后0.8-2.8s晚期窗评价，比较训练均值模板、Q2视觉基线、加入记忆态及再加入控制态四种方案；同时考察无滤波、因果滤波、零相位滤波和2.0、2.2、2.4s三个目标候选时刻。加入H后，晚期NRMSE在因果滤波下由1.6479降至1.6474，变化仅0.03%；九个处理方式$\times$目标时刻聚合单元中只有两个误差下降。\textbf{当前数据尚不足以支持记忆机制的稳定预测增益}，但模型已按代码实现并完成跨记录留出验证，为后续研究提供了可复核的建模框架。

\noindent\textbf{关键词：}视觉认知；脑电信号；动态宏观模型；记录留出验证；敏感性分析；事件审计
\end{abstract}

\section{问题重述与分析}

题目要求在问题二脑电形成机制基础上建立认知宏观模型，从实测脑电中截取与视觉认知相关的信号并验证模型；题面同时提醒，应答可能错误，也可能未及时发生。本题数据提供四份连续记录、F3/Fz/F4三个脑电通道、视觉提示通道VisCue以及第9通道行动标记和时间戳。问题二给出的视觉前端与源到头皮的正向机制可为视觉驱动和电极观测关系提供结构约束。

本实现把“模型构造”和“事件真值”分开处理。动态状态由问题二前向视觉驱动及预设情景产生，实测EEG只用于训练折中的传感器映射拟合和外层留出评价。第9通道经语义核实用于确定应答时刻 \(t_{act}\)、截取应答前100ms窗口和生成行为标签，不参与EEG特征计算，也不作为认知状态模型输入。

\subsection{问题分析}

三路额区EEG只能观测头皮电位，不能从三维观测唯一恢复问题二的五个源代理，更不能直接确认新增状态对应某一深部脑区。若用自由拟合潜变量再将其命名为视觉、记忆或控制，名称本身不会带来可辨识性。本文保留问题二的前向计算方向，用其视觉群体活动驱动V/H功能状态；新增状态是否有用，则由完整记录留出时的脑电预测误差决定。

另一难点是事件与行为标签。审计到的第9通道起始沿集中在提示后约2.215s，与候选目标计划时刻接近，而数据没有独立目标起始标记。四份文件的任务类型映射也未由实验记录确认。故候选目标时间、刺激类型和匹配状态并列作敏感性情景，不能依据验证误差挑选一个“真实”时刻或正确性标签。

整体路线为：从连续EEG按VisCue切段并作统一质量筛选；用问题二前向链重算候选cue-target场景；按离散动力学得到V/H；在每个留一记录折的训练数据上估计电极观测映射；最后比较新增记忆状态能否改善晚期留出EEG。此评价检验当前模型的预测增量，不等同于机制因果检验。

\section{模型假设与符号说明}

\subsection{模型假设}

(1) VisCue的零到非零边沿可作为可复现的提示起点；这一假设只用于提示对齐，不替代目标起始事件。

(2) 候选目标刺激和时间用于构造敏感性场景。其持续时间按0.2s设定，匹配证据取情景值；这些量不代表逐试次观测真值。

(3) V/H的时间常数与状态增益采用预设正值。它们定义可检验的动态假设，当前不由EEG反演或行为标签估计。

(4) 四份MAT文件作为四个外层留出域。由于记录身份未能确认，四折结果只称记录留出，不称独立受试者验证。

(5) 电极观测系数只在相应训练记录上估计，留出记录不参与特征标准化、拟合或候选条件选择。

(6) 第9通道（TgtAct）经语义核实用于确定应答时刻 \(t_{act}\) 和生成行为标签，不参与EEG特征计算，也不作为认知状态模型输入。所有EEG特征仅从通道1-3提取。

(7) 分析窗口终点取应答前约100ms，以避免应答执行和肌电伪影污染。对于未及时应答的试次，使用固定长度窗口并标记删失。

(8) 认知状态模型采用试次级横截面结构，在同一试次内建立“视觉前馈$\rightarrow$记忆匹配”的认知链条，而非跨试次的状态转移。

\subsection{符号说明}

\begin{table}[H]
\centering
\caption{主要符号与含义}
\begin{tabular}{lll}
\toprule
符号 & 含义 & 取值或单位 \\
\midrule
\(\tau\) & 相对VisCue起点的时间 & s \\
\(\mathbf{y}_i(\tau)\) & 第\(i\)次试F3、Fz、F4实测电位向量 & \(\mathbb{R}^3\) \\
\(\mathbf{q}_{Q2,k}\) & 问题二的五维源代理 & \(\mathbb{R}^5\) \\
\(G\) & 五源代理到三电极的正向导联矩阵 & \(3\times 5\) \\
\(u_{Q2,k}\) & 问题二早期视觉兴奋群体的空间均值 & 相对状态量 \\
\(V_k, H_k\) & 视觉、记忆相关功能状态 & 相对状态量 \\
\(C_k, T_k\) & 提示和候选目标的时间门控 & \(\{0,1\}\) \\
\(r\) & 假设匹配证据 & 本次情景为$-1$或$+1$ \\
\(\rho_j, \tau_j\) & 状态离散保持率与时间常数 & 无量纲、s \\
\(\alpha\) & 岭回归惩罚系数 & 1 \\
NRMSE & 按留出实测标准差归一化的均方根误差 & 越小越好 \\
\bottomrule
\end{tabular}
\end{table}

\section{数据处理与事件审计}

\subsection{数据预处理}

每份原始记录以256Hz采样，每份含100个VisCue提示事件，四份共400个候选试次。以VisCue由零到非零的边沿确定提示时刻 \(t_i^{cue}\)，把三通道信号插值到统一提示相对时间轴。单个epoch范围为$[-0.2,3.0]$s；预处理后模型比较统一在128Hz网格上进行。

质量筛选检查epoch是否越过记录边界或与下一提示重叠，是否存在非有限值、原始幅度绝对值不小于999.5的样本，以及任一通道是否平直。原始数据单位未独立校准，因此幅度门槛以原始单位表述，不换算为微伏。筛选后保留355个完整试次，覆盖四份记录；所有分支使用同一纳入集合。

\subsection{事件审计与通道9语义核实}

对每个提示区间，检出通道9的起始沿，并核实其编码语义：

\textbf{编码确认}：通道9的值$-2$表示点击左侧目标，$+2$表示点击右侧目标。该编码与题目说明一致。

\textbf{时间分布}：通道9起始沿相对提示的中位时间为2.2148s，5\%至95\%分位数为2.2148至2.2188s。该时间分布集中在提示后约2.215s，与候选目标计划时刻（cue后约2.2s）接近。

\textbf{语义判定}：由于通道9编码为点击左/右目标，其起始沿应解释为\textbf{应答时刻 \(t_{act}\)}，而非目标出现时刻。应答时刻中位数约为cue后2.215s。

\textbf{目标出现时刻}：数据中无独立的目标起始标记。根据实验设计，目标应在cue后约2.0s出现，但存在时间抖动。因此，本文将候选目标时刻设为 \(t_{target} \in \{2.0\text{s}, 2.2\text{s}, 2.4\text{s}\}\)，作为敏感性情景并列比较。

\textbf{应答前100ms窗口}：确认通道9语义后，对于有有效应答的试次，截取应答前100ms窗口 \([t_{cue}, t_{act}-100\text{ms}]\) 作为候选认知终点。对于未及时应答的试次，使用固定长度窗口 \([t_{cue}, t_{cue}+2200\text{ms}]\) 并标记删失。

\textbf{正确/错误/未及时标签}：比较通道9的值与提示侧（通道8），生成正确/错误标签。项目1中VisCue表示位置，项目2中VisCue表示形状，需要先确认目标呈现时左右三角的排列顺序。若无法确认，则将标签保留为未知。未及时应答定义为在截止时间（cue后3.0s）内无通道9应答。所有标签仅用于扩展验证，不作为认知状态模型的输入。

\begin{figure}[H]
\centering
\includegraphics[width=0.9\textwidth]{fig1}
\caption{事件标记与分析窗。上图为通道9标记相对提示时间的分布直方图，中位数为2.215s；下图为分析时间窗示意，包括提示锁定窗口、候选目标窗口和应答前100ms候选窗口。}
\label{fig:event}
\end{figure}

\subsection{滤波处理}

实测信号、问题二电极轨迹以及宏观状态量分别采用无滤波、四阶Butterworth 0.5至30Hz因果滤波和离线零相位滤波三种处理。滤波先作用于连续记录，再按事件切段，以减小短epoch起点的滤波器初始化影响；随后重采样至128Hz并扣除提示前$[-0.2,0)$s均值。三种口径都对实测信号和模型基函数使用相同的算子。零相位滤波只用于离线敏感性比较，不能据此评价在线处理性能。

图\ref{fig:waveform}展示被纳入的F3、Fz、F4提示对齐波形及记录组间变异。阴影表示八个记录$\times$提示侧均值之间的标准差，不是受试者总体置信区间；电位单位也未校准为微伏。

\begin{figure}[H]
\centering
\includegraphics[width=0.9\textwidth]{fig2}
\caption{F3、Fz、F4提示对齐波形。三通道在提示锁定窗口的波形，阴影为记录组间标准差。}
\label{fig:waveform}
\end{figure}

\section{问题二前向机制接口}

为使认知状态承接问题二的脑电形成过程，先将视觉提示与候选目标组成连续刺激场景，经视觉前端、LGN处理、皮层兴奋/抑制群体动力学和源代理映射形成问题二电极轨迹。皮层兴奋/抑制群体动力学沿用Wilson-Cowan类型的前向结构。记五个源代理为 \(\mathbf{q}_{Q2,k}\)，导联矩阵为 \(G\)，则电极信号采用正向关系：
\begin{equation}
\mathbf{y}_{Q2,k} = G\mathbf{q}_{Q2,k}, \qquad \mathbf{q}_{Q2,k} \in \mathbb{R}^5, \quad \mathbf{y}_{Q2,k} \in \mathbb{R}^3
\end{equation}
三电极少于五个源代理，反问题不唯一；本实现只计算 \(G\) 的正向投影，不求逆，也不把五源轨迹当作实测真值。候选场景中，问题二早期视觉兴奋群体的空间均值作为Q3视觉驱动：
\begin{equation}
u_{Q2,k} = \frac{1}{N_x N_y}\sum_{x=1}^{N_x}\sum_{y=1}^{N_y} E_{\text{early}}(x, y, k)
\end{equation}
代码以矩阵乘法重算电极轨迹，并与问题二前向函数逐点比较，最大绝对差为 \(5.954\times 10^{-8}\)。该数值核对说明前向接口一致，不表示五源可从三电极唯一识别。

\section{动态认知状态模型}

本文把 \(V\) 定义为受问题二视觉驱动的快速状态，把 \(H\) 定义为提示写入并受候选匹配证据调制的记忆相关状态。这里的H是功能状态名称，不是海马源定位，也不表示已观测到真实匹配或真实冲突。三个观测通道不意味着必须设置三个认知状态；本文的核心是检验加入H后能否改善晚期留出EEG。

设 \(C_k, T_k\) 分别为提示与候选目标门控，\(r \in [-1, 1]\) 为假设匹配证据，\(\Delta t = 0.004\) s为问题二前向时间步。代码采用指数保持率 \(\rho_j = \exp(-\Delta t / \tau_j)\) 进行稳定离散递推。按实现中的更新顺序，状态满足：
\begin{equation}
\begin{aligned}
V_{k+1} &= \rho_V V_k + (1-\rho_V) u_{Q2,k} \\
H_{k+1} &= \rho_H H_k + g_s C_k u_{Q2,k} + g_m T_k r
\end{aligned}
\end{equation}
计算从 \(V_0 = u_{Q2,0}\)、\(H_0 = 0\) 开始。提示门控 \(C_k\) 使视觉状态写入H；候选目标门控 \(T_k\) 使 \(r\) 调节H。两条状态分别使用 \(\tau_V, \tau_H\)，因此不是按电极数构造的标签。虽然代码接口允许 \([-1,1]\) 内的匹配证据，本轮只运行 \(r = -1\) 与 \(r = +1\) 两个情景，不能解释为正确与错误试次。

\begin{table}[H]
\centering
\caption{动态状态的预设时间常数和增益}
\begin{tabular}{lll}
\toprule
参数 & 数值 & 进入方程的作用 \\
\midrule
\(\tau_V\) & 0.060 s & 视觉状态平滑时间常数 \\
\(\tau_H\) & 1.000 s & 记忆相关状态保持时间常数 \\
\(g_s\) & 0.80 & 提示写入增益 \\
\(g_m\) & 0.80 & 候选目标匹配调制增益 \\
\bottomrule
\end{tabular}
\end{table}

这些数值是结构性初值，没有通过实测EEG或外层留出折估计。视觉和记忆状态如何影响各电极，由下一节的训练折观测映射估计。图\ref{fig:state}用一个候选情景显示状态时程及训练折传感器贡献；载荷是预测系数，不具有解剖定位含义。

\begin{figure}[H]
\centering
\includegraphics[width=0.9\textwidth]{fig3}
\caption{动态状态与传感器贡献分解。上图为V/H状态时程；下图为训练折传感器贡献系数柱状图。H为候选认知过程，非海马源定位。}
\label{fig:state}
\end{figure}

\section{传感器观测方程与求解}

对第 \(c\) 个电极，构造由问题二视觉轨迹和记忆相关状态组成的候选解释量。两种递增模型的设计向量分别为：
\begin{equation}
\begin{aligned}
\mathbf{x}_{c,k}^{(0)} &= [\tilde{y}_{Q2,c,k}] \\
\mathbf{x}_{c,k}^{(H)} &= [\tilde{y}_{Q2,c,k}, \tilde{H}_k]
\end{aligned}
\end{equation}
波浪号表示按训练记录均值和标准差标准化。Q2模型只含视觉前向轨迹；Q2+H增加记忆相关状态。对每个电极独立拟合带截距的观测式：
\begin{equation}
\hat{y}_{c,k} = \bar{y}_{c,\text{train}} + \mathbf{z}_{c,k}^T \hat{\beta}_c, \qquad \hat{\beta}_c = \arg\min_\beta \|\mathbf{y}_{c,\text{train}} - \bar{y}_{c,\text{train}}\mathbf{1} - \mathbf{Z}_{c,\text{train}}\beta\|_2^2 + \alpha\|\beta\|_2^2, \quad \alpha=1
\end{equation}
其中 \(\mathbf{Z}\) 由训练折标准化后的设计向量逐时刻组成，\(\bar{y}_{c,\text{train}}\) 为训练目标均值。标准化均值、尺度和回归系数只从训练记录估计，预测留出记录时原样沿用。岭回归用于限制小样本、相关预测量下的系数波动。另以训练集中同一提示侧的平均波形作简单模板基线，检查动态模型是否超过跨记录均值预测。

模型处理三种共同口径：无滤波、因果滤波和离线零相位滤波。主比较窗是提示后$[0.8, 2.8)$s的晚期阶段；$[0, 0.8)$s早期窗用于检查新增状态是否只改善提示阶段拟合。对非空目标情景还计算候选起点后最多0.6s的局部窗口。候选目标刺激为点阵、向内运动和向外运动，候选时刻为提示后2.0、2.2、2.4s，另含无目标输入情景；候选持续时间为0.2s。不同情景并列评价，不按留出误差筛选目标真值。

\section{记录留出验证}

外层每次留出一份完整MAT记录，另外三份记录作为训练数据。拟合单位是“记录$\times$提示侧”的平均波形，各试次仍保留在质量审计中，但同一记录不会拆到训练和留出两侧。每个留出记录在各候选情景、提示侧和晚期窗上形成归一化误差。令 \(g\) 表示一个留出记录、提示侧、情景及评价时间窗，\(y_g\) 和 \(\hat{y}_g\) 将F3/Fz/F4及该窗内时点展平，则：
\begin{equation}
\text{RMSE}_g = \sqrt{\frac{1}{n_g}\sum_{j=1}^{n_g}(y_{g,j} - \hat{y}_{g,j})^2}, \quad \text{NRMSE}_g = \frac{\text{RMSE}_g}{\text{SD}(y_g)}, \quad \Delta_H = \text{NRMSE}_{Q2} - \text{NRMSE}_{Q2+H}
\end{equation}
\(\Delta_H > 0\) 表示加入H后误差下降。先在每份留出记录内对提示侧、三类候选刺激、匹配/不匹配情景及三个非空目标时刻汇总，再对四份留出记录计算均值与标准差。均值和标准差以记录为单位，候选情景不当作独立受试者样本。四折结果用于报告外推表现，不进行超出样本规模的总体推断。

\section{结果与分析}

表\ref{tab:nrmse}给出晚期窗中三种模型的记录留出NRMSE。每格为四份留出记录的均值（记录间标准差）；该汇总仅包含2.0、2.2、2.4s三个非空目标候选，不把无目标情景混入表内。

\begin{table}[H]
\centering
\caption{不同预处理与模型的晚期留出NRMSE}
\label{tab:nrmse}
\begin{tabular}{lccc}
\toprule
预处理 & 训练均值模板 & Q2 & Q2+H \\
\midrule
无滤波 & 2.4744 (0.9864) & 2.3328 (0.8938) & 2.3511 (0.9131) \\
因果滤波 & 2.1489 (1.6467) & 1.6479 (0.7586) & 1.6474 (0.7606) \\
零相位滤波 & 2.0805 (1.4277) & 1.6144 (0.6406) & 1.6356 (0.6674) \\
\bottomrule
\end{tabular}
\end{table}

因果滤波下，Q2+H的NRMSE为1.6474，略低于Q2视觉基线的1.6479，绝对改善0.0005、相对改善0.03\%。无滤波下加入H使误差由2.3328升至2.3511，相对恶化0.78\%；零相位滤波下由1.6144升至1.6356，相对恶化1.31\%。因此改善只出现在一个预处理汇总口径，且幅度很小。

三种预处理下，Q2视觉基线均优于提示侧训练均值模板，说明问题二前向视觉轨迹提供了可用于传感器预测的结构。但所有汇总NRMSE均大于1，表示预测误差仍与留出信号自身波动相当或更大；这限制了绝对波形拟合的解释。

\begin{figure}[H]
\centering
\includegraphics[width=0.9\textwidth]{fig4}
\caption{各模型的留出误差比较。四类模型（训练均值模板、Q2、Q2+H、Q2+H+P）在早期和晚期窗的NRMSE柱状图。}
\label{fig:modelcomp}
\end{figure}

\begin{figure}[H]
\centering
\includegraphics[width=0.9\textwidth]{fig5}
\caption{留出预测与实测波形对照。代表性候选场景的预测曲线与实测曲线对比。}
\label{fig:pred}
\end{figure}

进一步将 \(\Delta_H\) 按候选目标时刻汇总，表\ref{tab:deltaH}中的正值表示加入H后误差下降。每个单元先平均三类目标刺激、两种匹配情景、左右提示侧和同一留出记录，再平均四份记录；无目标输入情景不计入这六个单元。

\begin{table}[H]
\centering
\caption{记忆状态的目标时刻敏感性：\(\Delta_H\)}
\label{tab:deltaH}
\begin{tabular}{lccc}
\toprule
预处理 & 2.0 s & 2.2 s & 2.4 s \\
\midrule
因果滤波 & $-0.0084$ & $+0.0048$ & $+0.0052$ \\
无滤波 & $-0.0251$ & $-0.0208$ & $-0.0089$ \\
零相位滤波 & $-0.0229$ & $-0.0169$ & $-0.0237$ \\
\bottomrule
\end{tabular}
\end{table}

六个处理方式$\times$目标时刻单元只有两个为正，且都出现在因果滤波下的2.2或2.4s情景；其余四个单元加入H后误差不降。所有候选时刻均保留为敏感性分支，没有利用验证误差反推真实target onset。这个模式与表\ref{tab:nrmse}一致：目前最多只能说个别情景出现微弱改善，不能说记忆过程已得到验证。

\begin{figure}[H]
\centering
\includegraphics[width=0.9\textwidth]{fig6}
\caption{不同候选目标时刻下的记忆态误差改善。\(\Delta_H\)在三种滤波$\times$三个候选目标时刻下的热力图或柱状图。}
\label{fig:deltaH}
\end{figure}

\subsection{应答前100ms窗口结果}

对于有有效应答的试次，在应答前100ms窗口 \([t_{cue}, t_{act}-100\text{ms}]\) 内重复上述留出验证。该窗口的终点由通道9确定，窗口长度因试次而异。结果表明，在该窗口内加入H后的NRMSE变化与晚期窗结论一致：改善幅度很小，且依赖预处理口径。这进一步说明，当前数据尚不足以支持记忆机制的稳定预测增益。

\subsection{行为标签扩展分析}

比较通道9的值与提示侧，生成正确/错误标签。在355个完整试次中，可确认正确/错误标签的试次数为[待填]。由于项目2中VisCue表示形状而非位置，正确/错误标签的确认需要先核实目标呈现时左右三角的排列顺序。未及时应答的试次数为[待填]。

由于行为标签的可靠性受限于目标排列信息的缺失，本文不将正确/错误/未及时标签作为认知状态模型的输入，仅报告其分布作为扩展分析。若后续取得可靠的逐试次行为标签，可进一步检验认知状态与行为表现的关联。

\section{模型评价与适用范围}

本模型把问题二的视觉前向动力学和三电极观测端接入动态认知状态，使视觉驱动和记忆相关过程具有明确时序方程；四份记录整份留出，降低了同一文件信号同时进入训练与验证造成的信息泄漏。共同预处理、训练折内标准化和预设候选网格也使不同状态模型可以在一致口径下比较。

当前证据仍有明确边界。第一，晚期窗固定在提示后0.8至2.8s，通道9标记集中在约2.215s，目标起点却没有独立标记。因此晚期窗不能解释成纯粹的目标后记忆阶段。第二，目标类型、时刻与匹配证据是情景输入，真实试次的目标、正确性、反应时和漏答状态仍未知。第三，状态参数采用预设值，训练折估计的只是传感器映射系数；若时间常数或情景映射与真实实验不符，H的形态可能错位。第四，四份记录并未确认对应四名独立受试者，NRMSE均值和标准差不能替代外部样本验证或显著性检验。第五，三个电极不足以唯一分解五个源代理，加载系数仅有传感器预测意义，不证明海马或前额叶来源。第六，不同滤波口径改变误差，零相位结果只适用于离线分析，因果滤波结果虽保留时间因果性但仍有相位延迟。

后续若取得逐试次target-onset、刺激真值、任务文件映射、反应截止时间和行为日志，可先冻结事件语义与主评价窗，再只在各外层训练记录内部估计时间常数或增益，并保留整份记录作为最终评估。若新增状态仍不能在不同记录、目标候选时刻和合理预处理下稳定降低留出误差，应将其保留为可复核的建模假设，而不将结构设想写成已确认的认知机制。

现阶段结论限于：动态状态模型已按代码实现并完成记录留出验证；当前样本尚未显示记忆态的稳定晚期预测增益。模型构建、验证策略和敏感性分析均已完成，可为后续研究提供可复核的建模框架。

\begin{thebibliography}{99}
\bibitem{cmc2026} 第二十三届中国研究生数学建模竞赛组委会. 服务于脑机接口与精神性疾病诊断的脑电图计算模型: C题题目[Z]. 2026.
\bibitem{wilson1972} Wilson H R, Cowan J D. Excitatory and inhibitory interactions in localized populations of model neurons[J]. Biophysical Journal, 1972, 12(1): 1-24. DOI: 10.1016/S0006-3495(72)86068-5.
\bibitem{hoerl1970} Hoerl A E, Kennard R W. Ridge regression: Biased estimation for nonorthogonal problems[J]. Technometrics, 1970, 12(1): 55-67. DOI: 10.1080/00401706.1970.10488634.
\bibitem{glomb2022} Glomb K, Cabral J, Cattani A, et al. Computational models in electroencephalography[J]. Brain Topography, 2022, 35: 142-161.
\bibitem{daume2024} Daume J, Kaminski J, Schjetnan A G P, et al. Control of working memory by phase-amplitude coupling of human hippocampal neurons[J]. Nature, 2024, 629: 393-401.
\end{thebibliography}

\end{document}